%% elimina il numero di pagina della corrente pagina:
\thispagestyle{empty}
%% fa si' che la pagina corrente venga contata come pagina 0, 
%% cosi' la prossima pagina comincia da 1.
\setcounter{page}{0}

\begin{center}
\Large
\textsc{Università degli Studi di Torino}\\
\vspace{-0.4cm}
\rule{\textwidth}{0.1mm}\\
\large
\textsc{Dipartimento di Informatica}\\ 
\bigskip
Corso di Laurea in Informatica \\
\vfill
\includegraphics[scale=0.5]{universita}
\vfill %%\vspace{3cm}
Tesi di Laurea in Informatica \\
\vfill %%\vspace{3cm}
\Large
\textsc{BLINC}\\
\textsc{Blockchain e Identità Digitale}
\end{center}
\vfill
%%%%%%
%%%%%% Qui va inserito il nome del/lla laureando/a
\begin{tabular}[t]{l}
Candidato:\\
Lorenzo Bersano
\end{tabular}
%%%%%%
\hfill 
%%%%%% Qui va inserito il nome del relatore
\begin{tabular}[t]{l}
Relatori: \\
Prof. Claudio Schifanella \\
Prof. Alex Cordero
\end{tabular}
%%%%%%
\vfill
\begin{center}
\normalsize
\rule{8cm}{0.1mm}\\
\bigskip
ANNO ACCADEMICO 2017--2018
\end{center}
